\section{Definition of abelian varieties}
\label{mumford-II4}
\dcref{M:ch2:II.4}

\begin{convention}[Varieties and completeness]
\label{mumford-conv-variety-field}
\dcref{M:ch2:II.4}
The ground field $k$ is algebraically closed.  In this chapter a variety is an
irreducible separated scheme of finite type over $k$; complete means proper
over $k$.  Products and morphisms are taken over $k$.
\notready
\end{convention}

\begin{definition}[Abelian variety]
\label{mumford-def-abelian-variety}
\dcref{M:ch2:II.4}
An abelian variety is a complete variety $A$ equipped with a point $0$, a
multiplication $m:A\times A\to A$, and an inverse $i:A\to A$ which are
morphisms and satisfy the group axioms.  We write $x+y=m(x,y)$ and
$-x=i(x)$.
\uses{mumford-conv-variety-field}
\notready
\end{definition}

\begin{lemma}[Translations]
\label{mumford-lem-abelian-translations}
\dcref{M:ch2:II.4}
For every $a\in A$, the translation $t_a(x)=x+a$ is an automorphism of $A$;
its inverse is $t_{-a}$.  Consequently all local geometric properties of $A$
are invariant under translation.
\uses{mumford-def-abelian-variety}
\begin{proof}
The group identities give $t_{-a}t_a=t_at_{-a}=\operatorname{id}$.  Both maps
are morphisms because they are composites of the multiplication and inverse
morphisms.
\end{proof}
\notready
\end{lemma}

\begin{proposition}[Multiplication by an integer]
\label{mumford-prop-multiplication-n}
\dcref{M:ch2:II.4}
For $n\in\Z$, let $[n]:A\to A$ be repeated addition (and inversion for
$n<0$).  If $n$ is nonzero and prime to $\operatorname{char}(k)$, then $[n]$
is finite and etale, hence surjective.  Its kernel $A[n]$ is a finite group
scheme of rank $n^{2\dim A}$.
\uses{mumford-def-abelian-variety,mumford-lem-abelian-translations}
\notready
\end{proposition}

\begin{lemma}[Rigidity Lemma, Form I]
\label{mumford-lem-rigidity-form-I}
Let $X$ be complete and let $Y,Z$ be varieties.  If
$f:X\times Y\to Z$ is a morphism and there are $y_0\in Y$, $z_0\in Z$ with
$f(X\times\{y_0\})=\{z_0\}$, then there is a unique morphism $g:Y\to Z$ such
that $f=g\circ p_2$.
\dcref{M:ch2:II.4}
\uses{mumford-conv-variety-field}
\begin{proof}
Restrict $f$ to the slice through a chosen $x_0$; this gives the only possible
$g$.  To verify that it works, test the graph of $f$ over an affine
neighbourhood of $z_0$.  Properness of the first projection shows that the
parameters for which the image leaves this neighbourhood form a closed set
disjoint from $y_0$.  Over its complement, regular functions on the complete
fibres are constant, so the graph is the pullback of the graph of $g$.
Equality of morphisms on this dense open and separatedness of $Z$ extend the
identity to all of $X\times Y$; the slice also proves uniqueness.
\end{proof}
\notready
\end{lemma}

\begin{corollary}[Morphisms of abelian varieties]
\label{mumford-cor-morphism-translation}
If $A,B$ are abelian varieties and $f:A\to B$ is a morphism, then
$f(x)=h(x)+b$ for a unique $b\in B$ and a homomorphism $h:A\to B$.
\dcref{M:ch2:II.4}
\uses{mumford-lem-rigidity-form-I,mumford-def-abelian-variety}
\begin{proof}
Translate the target so that $f(0)=0$.  Compare $f\circ m$ with the sum of
the two pullbacks of $f$.  Their quotient is a morphism from $A\times A$ to
$B$ which is trivial on both coordinate axes.  Applying rigidity first in one
variable and then in the other makes this quotient trivial, proving additivity.
The original translation is the unique constant term.
\end{proof}
\notready
\end{corollary}

\begin{corollary}[Commutativity]
\label{mumford-cor-abelian-commutative}
The abstract group underlying an abelian variety is commutative.
\dcref{M:ch2:II.4}
\uses{mumford-cor-morphism-translation}
\begin{proof}
Apply the morphism decomposition to inversion after translating its value at
the identity (which is already the identity).  Comparing inversion of a sum
with the two possible orders gives $x^{-1}+y^{-1}=y^{-1}+x^{-1}$; invert once
more to obtain commutativity.
\end{proof}
\notready
\end{corollary}

\begin{corollary}[Pointed Hom is linear]
\label{mumford-cor-pointed-hom-linear}
For complete pointed varieties $S,T$ and an abelian variety $A$, restriction to
the two coordinate axes gives a bijection
\[
 \Hom_*(S,A)\times\Hom_*(T,A)\xrightarrow{\sim}
 \Hom_*(S\times T,A),\qquad (f,g)\mapsto((s,t)\mapsto f(s)+g(t)).
\]
\dcref{M:ch2:II.4}
\uses{mumford-lem-rigidity-form-I,mumford-cor-abelian-commutative}
\begin{proof}
Normalize $h$ by removing its two axis restrictions.  The normalized map is
zero on each axis, and rigidity in either variable forces it to be zero on the
whole product.  The two restrictions therefore reconstruct $h$ uniquely, and
the group law supplies the displayed sum.
\end{proof}
\notready
\end{corollary}

\begin{theorem}[Unital composition criterion]
\label{mumford-thm-unital-composition}
Let $X$ be complete, let $e\in X$, and let $m:X\times X\to X$ be a morphism
with $m(x,e)=m(e,x)=x$.  Then associativity and the existence of inverse follow
from completeness and rigidity, so $(X,m,e)$ is an abelian variety.
\dcref{M:ch2:II.4}
\uses{mumford-lem-rigidity-form-I}
\notready
\end{theorem}

\begin{corollary}[No rational curves]
\label{mumford-cor-no-rational-curves}
Every morphism $\PP^1_k\to A$ to an abelian variety is constant.  The same
holds for a unirational complete variety.  This is a derived corollary of the
rigidity lemma.
\dcref{M:ch2:II.4}
\uses{mumford-def-abelian-variety,mumford-lem-rigidity-form-I}
\notready
\end{corollary}

\section{Supporting dependency interfaces}
\label{mumford-II-coverage-register}
These compact nodes record auxiliary results and applications used by the
variety-theoretic constructions.  Detailed proofs are treated as external
inputs where the available evidence is insufficient for reconstruction.

\begin{remark}[First commutativity argument]
\label{ch02-sec04-commutativity-first-proof}
The invariant-differential argument gives commutativity independently of the
translation decomposition; it applies when the group law is transported
through a quotient.
\dcref{M:ch2:II.4}
\uses{mumford-def-abelian-variety}
\notready
\end{remark}

\begin{remark}[Cohomology fibre notation]
\label{ch02-sec05-fiber-notation}
For $y\in Y$, write $X_y=X\times_Y\Spec k(y)$ and
$\cF_y=\cF\otimes_{\cO_Y}k(y)$.  All semicontinuity and base-change claims
use this restricted sheaf notation.
\dcref{M:ch2:II.5}
\uses{mumford-def-f-flat}
\notready
\end{remark}

\begin{lemma}[Constant-fibre freeness]
\label{ch02-sec05-constant-fiber-lemma}
If the fibre dimension of a cohomology group is locally constant on a reduced
base, the corresponding cohomology sheaf is locally free and commutes with
base change.
\dcref{M:ch2:II.5}
\uses{mumford-thm-base-change-criterion}
\notready
\end{lemma}

\begin{lemma}[Split image in a finite complex]
\label{ch02-sec05-image-splitting-lemma}
On a locus where the rank of a differential in the finite projective complex is
constant, its image is locally a direct summand; this is the linear-algebra
step behind the base-change criterion.
\dcref{M:ch2:II.5}
\uses{mumford-thm-grothendieck-complex}
\notready
\end{lemma}

\begin{lemma}[Trivial line-bundle section test]
\label{ch02-sec05-trivial-line-bundle-criterion}
On a complete connected variety, a line bundle is trivial exactly when it and
its inverse have nonzero global sections whose product is a nonzero constant.
\dcref{M:ch2:II.5}
\uses{mumford-thm-proper-direct-image-coherent}
\notready
\end{lemma}

\begin{remark}[Cube applications]
\label{ch02-sec06-unread-proof}
Further symmetry and Picard-group identities follow from the cube; these
applications are external inputs while the displayed cube and
cross-effect formulation provide their dependencies.
\dcref{M:ch2:II.6}
\uses{mumford-thm-cube-varieties,mumford-prop-picard-quadratic}
\notready
\end{remark}

\begin{remark}[Finite-quotient coverage]
\label{ch02-sec07-coverage}
The finite-quotient construction also treats orbit sheaves, principal bundles,
and translation quotients; the finite-quotient theorem and translation-isogeny
proposition are the interfaces used by subsequent arguments.
\dcref{M:ch2:II.7}
\uses{mumford-thm-finite-group-quotient,mumford-prop-translation-quotient-isogeny}
\notready
\end{remark}

\begin{remark}[Remaining quotient applications]
\label{ch02-sec07-unread-results}
The affine-local quotient calculations and their descent proofs are retained as
external inputs.
\dcref{M:ch2:II.7}
\uses{ch02-sec07-coverage}
\notready
\end{remark}

\begin{remark}[Characteristic-zero dual construction]
\label{ch02-sec08-coverage}
The characteristic-zero Picard construction, normalized Poincare bundle,
biduality, and functorial dual maps form one dependency package.
\dcref{M:ch2:II.8}
\uses{mumford-thm-dual-char-zero,mumford-prop-biduality-char-zero}
\notready
\end{remark}

\begin{remark}[Remaining duality arguments]
\label{ch02-sec08-unread-results}
The representability and Poincare descent calculations not reproduced in detail
are external inputs for the scheme-theoretic duality construction.
\dcref{M:ch2:II.8}
\uses{ch02-sec08-coverage}
\notready
\end{remark}

\begin{remark}[Complex comparison coverage]
\label{ch02-sec09-coverage}
The analytic comparison identifies the dual lattice and Poincare factor of
automorphy; its remaining period-matrix computations are external inputs.
\dcref{M:ch2:II.9}
\uses{mumford-thm-analytic-dual-comparison}
\notready
\end{remark}

\begin{remark}[Complex-comparison proof interface]
\label{ch02-sec09-unread-results}
The omitted examples and detailed analytic identifications are external inputs;
the comparison theorem supplies the interface needed by the complex case.
\dcref{M:ch2:II.9}
\uses{ch02-sec09-coverage}
\notready
\end{remark}

\begin{remark}[Abelian-variety structure package]
\label{ch02-sec04-abelian-variety}
The complete group-variety definition is the standing object for this theory.
\dcref{M:ch2:II.4}
\uses{mumford-def-abelian-variety}
\notready
\end{remark}

\begin{remark}[Rigidity package]
\label{ch02-sec04-rigidity}
The Form I rigidity lemma is the primitive from which the morphism and
no-rational-curves consequences are derived.
\dcref{M:ch2:II.4}
\uses{mumford-lem-rigidity-form-I}
\notready
\end{remark}

\begin{remark}[Translation decomposition package]
\label{ch02-sec04-morphism-translation}
Every morphism between abelian varieties is a unique translation followed by a
homomorphism.
\dcref{M:ch2:II.4}
\uses{mumford-cor-morphism-translation}
\notready
\end{remark}

\begin{remark}[Second commutativity argument]
\label{ch02-sec04-commutativity-second-proof}
The inverse-morphism argument gives a second proof of commutativity in the
abelian-variety setting.
\dcref{M:ch2:II.4}
\uses{mumford-cor-abelian-commutative}
\notready
\end{remark}

\begin{remark}[Product-linearity package]
\label{ch02-sec04-hom-linearity}
Pointed maps from a product into an abelian variety split as the sum of their
two axis restrictions.
\dcref{M:ch2:II.4}
\uses{mumford-cor-pointed-hom-linear}
\notready
\end{remark}

\begin{remark}[Unital-law package]
\label{ch02-sec04-unital-law}
A unital multiplication on a complete variety satisfies the remaining group
axioms by rigidity.
\dcref{M:ch2:II.4}
\uses{mumford-thm-unital-composition}
\notready
\end{remark}

\begin{remark}[Projective-line package]
\label{ch02-sec04-no-rational-curves}
Maps from $\PP^1$ and, more generally, unirational complete varieties to an
abelian variety are constant.
\dcref{M:ch2:II.4}
\uses{mumford-cor-no-rational-curves}
\notready
\end{remark}

\begin{remark}[Proper direct-image package]
\label{ch02-sec05-proper-direct-image}
Proper pushforward preserves coherence for higher direct images of coherent
sheaves.
\dcref{M:ch2:II.5}
\uses{mumford-thm-proper-direct-image-coherent}
\notready
\end{remark}

\begin{remark}[Flatness package]
\label{ch02-sec05-flatness}
Flatness over the base is checked stalkwise for the structure morphism.
\dcref{M:ch2:II.5}
\uses{mumford-def-f-flat}
\notready
\end{remark}

\begin{remark}[Finite-complex package]
\label{ch02-sec05-finite-projective-complex}
A finite projective complex computes all fibre cohomology after arbitrary base
change in a proper flat family.
\dcref{M:ch2:II.5}
\uses{mumford-thm-grothendieck-complex}
\notready
\end{remark}

\begin{remark}[Complex-model package]
\label{ch02-sec05-complex-model}
The bounded complex can be replaced by a flat finite model without changing
cohomology.
\dcref{M:ch2:II.5}
\uses{mumford-lem-flat-complex-replacement}
\notready
\end{remark}

\begin{remark}[Complex base-change package]
\label{ch02-sec05-complex-base-change}
Quasi-isomorphisms of bounded flat complexes remain quasi-isomorphisms after
tensoring with a base algebra.
\dcref{M:ch2:II.5}
\uses{mumford-lem-flat-quasi-iso-base-change}
\notready
\end{remark}

\begin{remark}[Semicontinuity package]
\label{ch02-sec05-semicontinuity}
Fibre cohomology dimensions are upper semicontinuous and Euler characteristic is
locally constant.
\dcref{M:ch2:II.5}
\uses{mumford-thm-semicontinuity}
\notready
\end{remark}

\begin{remark}[Base-change criterion package]
\label{ch02-sec05-base-change-criterion}
Constant fibre dimension on a reduced connected base is equivalent to local
freeness and the cohomology-and-base-change isomorphism.
\dcref{M:ch2:II.5}
\uses{mumford-thm-base-change-criterion}
\notready
\end{remark}

\begin{remark}[Vanishing in the previous degree]
\label{ch02-sec05-vanishing-previous-degree}
Vanishing of fibre cohomology in degree $p$ forces base change in degree $p-1$.
\dcref{M:ch2:II.5}
\uses{mumford-cor-flat-base-change}
\notready
\end{remark}

\begin{remark}[Higher-vanishing package]
\label{ch02-sec05-higher-vanishing}
Vanishing of higher direct images bounds the degrees in which fibre cohomology
can occur.
\dcref{M:ch2:II.5}
\uses{mumford-cor-flat-base-change}
\notready
\end{remark}

\begin{remark}[Flat base-change package]
\label{ch02-sec05-flat-base-change}
Cohomology commutes with a flat extension of the base ring.
\dcref{M:ch2:II.5}
\uses{mumford-cor-flat-base-change}
\notready
\end{remark}

\begin{remark}[Seesaw package]
\label{ch02-sec05-seesaw}
The scheme of trivial fibres carries the residual line bundle whose pullback is
the original family.
\dcref{M:ch2:II.5}
\uses{mumford-thm-seesaw}
\notready
\end{remark}

\begin{remark}[Cube package]
\label{ch02-sec06-cube-theorem}
Three trivial coordinate faces force a line bundle on a triple product to be
trivial.
\dcref{M:ch2:II.6}
\uses{mumford-thm-cube-varieties}
\notready
\end{remark}

\begin{remark}[Product-splitting package]
\label{ch02-sec06-product-splitting}
The alpha--beta maps split a contravariant functor on a product into coordinate
and cross-effect parts.
\dcref{M:ch2:II.6}
\uses{mumford-prop-picard-quadratic}
\notready
\end{remark}

\begin{remark}[Order-of-functor package]
\label{ch02-sec06-order-functor}
The Picard functor has order at most two in the cross-effect filtration.
\dcref{M:ch2:II.6}
\uses{mumford-prop-picard-quadratic}
\notready
\end{remark}

\begin{remark}[Finite quotient package]
\label{ch02-sec07-quotient}
Finite group actions admit a finite quotient, and free actions give principal
etale/flat torsors according to the characteristic.
\dcref{M:ch2:II.7}
\uses{mumford-thm-finite-group-quotient}
\notready
\end{remark}

\begin{remark}[Characteristic-zero duality package]
\label{ch02-sec08-dual}
In characteristic zero, $\Pic^0(A)$ is represented by an abelian variety with
normalized Poincare bundle and biduality.
\dcref{M:ch2:II.8}
\uses{mumford-thm-dual-char-zero}
\notready
\end{remark}

\begin{remark}[Complex comparison package]
\label{ch02-sec09-comparison}
The analytic dual lattice and the algebraic Poincare bundle agree under
complexification.
\dcref{M:ch2:II.9}
\uses{mumford-thm-analytic-dual-comparison}
\notready
\end{remark}

\section{Cohomology and base change}
\label{mumford-II5}
\dcref{M:ch2:II.5}

\begin{definition}[Flatness over a base]
\label{mumford-def-f-flat}
\dcref{M:ch2:II.5}
For a morphism $f:X\to Y$ and a quasi-coherent sheaf $\cF$ on $X$, say that
$\cF$ is $f$-flat if each stalk $\cF_x$ is flat over $\cO_{Y,f(x)}$.
\notready
\end{definition}

\begin{theorem}[Coherence of proper direct images]
\label{mumford-thm-proper-direct-image-coherent}
If $f:X\to Y$ is proper between locally noetherian schemes and $\cF$ is
coherent, then each $R^pf_*\cF$ is coherent.
\dcref{M:ch2:II.5}
\notready
\end{theorem}

\begin{theorem}[Finite complex for a proper flat family]
\label{mumford-thm-grothendieck-complex}
Let $Y=\Spec A$, $f:X\to Y$ proper noetherian, and $\cF$ coherent and
$f$-flat.  There is a bounded complex $K^\bullet$ of finitely generated
projective $A$-modules such that, for every $A$-algebra $B$,
\[
H^p(X\times_Y\Spec B,\cF\otimes_A B)\simeq
H^p(K^\bullet\otimes_A B)
\]
functorially in $B$.
\dcref{M:ch2:II.5.1}
\uses{mumford-def-f-flat,mumford-thm-proper-direct-image-coherent}
\notready
\end{theorem}

\begin{lemma}[Flat replacement of a finite complex]
\label{mumford-lem-flat-complex-replacement}
Let $C^\bullet$ be a bounded complex over a noetherian ring with finitely
generated cohomology.  It admits a quasi-isomorphic complex $K^\bullet$ whose
positive terms are finite free and whose degree-zero term is flat whenever all
$C^p$ are flat.
\dcref{M:ch2:II.5}
\uses{mumford-thm-proper-direct-image-coherent}
\notready
\end{lemma}

\begin{lemma}[Flat quasi-isomorphisms survive base change]
\label{mumford-lem-flat-quasi-iso-base-change}
A quasi-isomorphism between bounded complexes of flat $A$-modules remains a
quasi-isomorphism after tensoring with every $A$-algebra.
\dcref{M:ch2:II.5}
\uses{mumford-lem-flat-complex-replacement}
\notready
\end{lemma}

\begin{theorem}[Semicontinuity and Euler characteristic]
\label{mumford-thm-semicontinuity}
For a proper morphism $f:X\to Y$ of noetherian schemes and a coherent sheaf
flat over $Y$, the function
$y\mapsto\dim_{k(y)}H^p(X_y,\cF_y)$ is upper semicontinuous, while
$\chi(\cF_y)=\sum_p(-1)^p\dim_{k(y)}H^p(X_y,\cF_y)$ is locally constant.
\dcref{M:ch2:II.5}
\uses{mumford-thm-grothendieck-complex,mumford-lem-flat-quasi-iso-base-change}
\notready
\end{theorem}

\begin{theorem}[Cohomology-and-base-change criterion]
\label{mumford-thm-base-change-criterion}
Assume $Y$ is reduced and connected, and let $K^\bullet$ be the finite complex
representing the derived pushforward.  Constancy of $\dim H^p(X_y,\cF_y)$ is equivalent to local freeness of
$R^pf_*\cF$ together with the fibrewise isomorphism
\[
 (R^pf_*\cF)\otimes k(y)\xrightarrow{\sim}H^p(X_y,\cF_y).
\]
Under these conditions the analogous map in degree $p-1$ is also an
isomorphism.
\dcref{M:ch2:II.5}
\uses{mumford-thm-grothendieck-complex,mumford-thm-semicontinuity}
\notready
\end{theorem}

\begin{corollary}[Vanishing and flat base change]
\label{mumford-cor-flat-base-change}
If $H^p(X_y,\cF_y)=0$ for every $y$, then base change in degree $p-1$ is an
isomorphism even when $Y$ is nonreduced.  If $R^kf_*\cF=0$ for $k>k_0$, then
$H^k(X_y,\cF_y)=0$ for $k>k_0$.  Finally, for a flat $A$-algebra $B$,
\[
 H^p(X\times_Y\Spec B,\cF\otimes_A B)\simeq H^p(X,\cF)\otimes_A B.
\]
\dcref{M:ch2:II.5}
\uses{mumford-thm-base-change-criterion,mumford-lem-flat-quasi-iso-base-change}
\notready
\end{corollary}

\begin{theorem}[Seesaw theorem]
\label{mumford-thm-seesaw}
Let $X$ be complete and connected, $T$ a noetherian variety, and $L$ a line
bundle on $X\times T$.  The scheme-theoretic locus
\[
 T_1=\{t\in T\mid L|_{X\times\{t\}}\simeq\cO_X\}
\]
is closed, and there is a line bundle $M$ on $T_1$ with
$L|_{X\times T_1}\simeq p_2^*M$.
\dcref{M:ch2:II.5}
\uses{mumford-thm-semicontinuity,mumford-cor-flat-base-change}
\begin{proof}
Apply the finite-complex description to $L$ and $L^{-1}$.  The simultaneous
rank-one locus for their zeroth cohomology sheaves is represented by a closed
subscheme $T_1$.  On this locus cohomology and base change make
$p_{2*}L$ invertible.  The universal evaluation map from its pullback to $L$
is an isomorphism on every geometric fibre, hence is an isomorphism by the
fibre criterion for coherent maps.  Taking $M=p_{2*}L$ gives the claim.
\end{proof}
\notready
\end{theorem}

\section{The theorem of the cube}
\label{mumford-II6}
\dcref{M:ch2:II.6}

\begin{theorem}[Theorem of the cube for varieties]
\label{mumford-thm-cube-varieties}
Let $X,Y$ be complete varieties, let $Z$ be a variety, and choose points
$x_0,y_0,z_0$.  If a line bundle $L$ on $X\times Y\times Z$ restricts
trivially to each coordinate face
$\{x_0\}\times Y\times Z$, $X\times\{y_0\}\times Z$, and
$X\times Y\times\{z_0\}$, then $L$ is trivial.
\dcref{M:ch2:II.6}
\uses{mumford-thm-seesaw}
\notready
\end{theorem}

\begin{proposition}[Picard splitting and order]
\label{mumford-prop-picard-quadratic}
For a contravariant functor $T$ from complete varieties to abelian groups, use
alternating sums of pullbacks along coordinate faces to form the cross-effects.
The projection and section maps then give a canonical splitting
\[
 T(X_0\times\cdots\times X_n)=\operatorname{im}(\alpha_T^n)\oplus
 \ker(\beta_T^n).
\]
The functor has order at most $n$ when its $(n+1)$-fold cross-effect vanishes.
For $T=\Pic$, the cube theorem is precisely the assertion that this order is
at most two.
\dcref{M:ch2:II.6}
\uses{mumford-thm-cube-varieties}
\notready
\end{proposition}

\section{Dividing varieties by finite groups}
\label{mumford-II7}
\dcref{M:ch2:II.7}

\begin{definition}[Free finite group action]
\label{mumford-def-free-finite-action}
\dcref{M:ch2:II.7}
A finite group $G$ acts freely on a variety $X$ if the map
$G\times X\to X\times X$, $(g,x)\mapsto(gx,x)$ is a closed immersion on
each component and distinct geometric points have distinct translates.
\notready
\end{definition}

\begin{theorem}[Finite quotient]
\label{mumford-thm-finite-group-quotient}
If a finite group acts on a quasi-projective variety $X$, the quotient sheaf
$X/G$ is represented by a variety and the quotient map is finite.  If the
action is free, the map is etale and $X\to X/G$ is a principal $G$-bundle in
the etale topology.
\dcref{M:ch2:II.7}
\uses{mumford-def-free-finite-action,mumford-def-abelian-variety}
\notready
\end{theorem}

\begin{proposition}[Translation quotients and isogenies]
\label{mumford-prop-translation-quotient-isogeny}
If $H\subset A$ is a finite subgroup acting by translations, then $A/H$ is an
abelian variety and the quotient map is an isogeny.  Conversely, every finite
etale isogeny is obtained this way from its kernel.
\dcref{M:ch2:II.7}
\uses{mumford-thm-finite-group-quotient,mumford-prop-multiplication-n}
\notready
\end{proposition}

\section{The dual abelian variety in characteristic zero}
\label{mumford-II8}
\dcref{M:ch2:II.8}

\begin{definition}[Algebraically trivial line bundles]
\label{mumford-def-piczero}
\dcref{M:ch2:II.8}
For a complete variety $X$, write $\Pic^0(X)$ for the connected component of
the identity in the Picard scheme, equivalently the subgroup of line bundles
algebraically equivalent to zero when the Picard scheme is represented.
\notready
\end{definition}

\begin{theorem}[Dual abelian variety in characteristic zero]
\label{mumford-thm-dual-char-zero}
For an abelian variety $A$ over a field of characteristic zero, $\Pic^0(A)$ is
represented by an abelian variety $\Adual$.  There is a normalized Poincare line
bundle $\cP$ on $A\times\Adual$, trivial on both coordinate origins, and it
represents the universal family of algebraically trivial line bundles on $A$.
\dcref{M:ch2:II.8}
\uses{mumford-def-piczero,mumford-thm-cube-varieties,mumford-thm-seesaw}
\notready
\end{theorem}

\begin{proposition}[Biduality and functoriality]
\label{mumford-prop-biduality-char-zero}
The Poincare bundle induces a canonical isomorphism $A\xrightarrow{\sim}
\widehat{\Adual}$ and reverses arrows: a homomorphism $f:A\to B$ has a
dual $\widehat f:\Bdual\to\Adual$ characterized by
$(f\times1)^*\cP_B\simeq(1\times\widehat f)^*\cP_A$.
\dcref{M:ch2:II.8}
\uses{mumford-thm-dual-char-zero}
\notready
\end{proposition}

\section{The case of the complex ground field}
\label{mumford-II9}
\dcref{M:ch2:II.9}

\begin{theorem}[Analytic description of the dual]
\label{mumford-thm-analytic-dual-comparison}
If $A(\C)=V/\Lambda$ is the analytic torus attached to a complex abelian
variety, then $\Adual(\C)$ is the torus whose lattice is dual to $\Lambda$ under
the integral pairing supplied by the Poincare bundle.  The analytic Poincare
bundle agrees with the Appell--Humbert bundle of that pairing.
\dcref{M:ch2:II.9}
\uses{mumford-cor-complex-av-torus,mumford-prop-biduality-char-zero,mumford-thm-appell-humbert}
\notready
\end{theorem}
